% C3. Chapinero es más caro y tiene peor información
\begin{frame}
  \begin{columns}[c]
    \begin{column}{0.6\textwidth}\centering
      \mapavertical{mapa_02_precio_m2_upz.pdf}
    \end{column}
    \begin{column}{0.38\textwidth}
      \titulolateral{Chapinero es más caro, casi todo apartamentos y con ubicación más imprecisa}
      \resizebox{\linewidth}{!}{%
      \begin{tabular}{@{}lrr@{}}
        \toprule
        & Chapinero & Bogotá \\
        \midrule
        Precio mediano (M)$^{*}$ & \preciomdchap & \preciomdresto \\
        Apartamentos & \aptotest & \aptotrain \\
        Coordenadas redondeadas & \coordredondeada & \coordredtrain \\
        \bottomrule
      \end{tabular}}

      \vspace{0.6em}
      {\scriptsize\color{gris}Apartamentos y coordenadas: test (Chapinero) frente a train (Bogotá). $^{*}$ Test no tiene precio: se compara Chapinero con el resto de Bogotá, ambos en train.\par}

      \vspace{0.6em}
      {\scriptsize\color{gris}Coordenadas redondeadas: con 3 decimales o menos, un error de ubicación de unos 110 m.\par}
    \end{column}
  \end{columns}
  \note{Chapinero está entre las zonas más caras de la ciudad. En test casi todo son apartamentos, y la ubicación es menos precisa: más de la mitad de las coordenadas viene redondeada. Además, al \sinareatest{} de test le falta el área, frente a solo el \sinareachap{} de los anuncios de Chapinero que sí están en train.}
\end{frame}
